\documentclass[11pt,a4paper]{article}
\usepackage[T1]{fontenc}
\usepackage{lmodern}
\usepackage[margin=27mm,headheight=14pt]{geometry}
\usepackage{amsmath,amssymb,amsthm,mathtools}
\usepackage{microtype,booktabs,tabularx,array,enumitem}
\usepackage{xcolor,fancyhdr}
\usepackage[unicode,colorlinks=true,linkcolor=blue!40!black,citecolor=blue!40!black,urlcolor=blue!40!black]{hyperref}
\hypersetup{pdftitle={Smooth isometric embeddings of closed orientable surfaces in four-space},pdfauthor={},pdfsubject={Smooth isometric embeddings of closed orientable surfaces},pdfkeywords={isometric embedding, closed surface, four-space, transverse oscillation}}
\newcommand{\R}{\mathbb R}
\newcommand{\Sph}{\mathbb S}
\newcommand{\dd}{\mathrm d}
\newcommand{\Hess}{\operatorname{Hess}}
\newcommand{\Span}{\operatorname{span}}
\newcommand{\avg}{\operatorname{avg}}
\newtheorem{theorem}{Theorem}[section]
\newtheorem{proposition}[theorem]{Proposition}
\newtheorem{lemma}[theorem]{Lemma}
\theoremstyle{definition}
\newtheorem{definition}[theorem]{Definition}
\theoremstyle{remark}
\newtheorem{remark}[theorem]{Remark}
\numberwithin{equation}{section}
\setlength{\parindent}{1.15em}
\setlength{\parskip}{3pt}
\setlist{itemsep=2pt,topsep=4pt,leftmargin=*}
\emergencystretch=1.3em
\pagestyle{fancy}
\fancyhf{}
\fancyhead[L]{\small Smooth isometric embeddings in four-space}
\fancyhead[R]{\small Public draft 1}
\fancyfoot[C]{\small\thepage}
\renewcommand{\headrulewidth}{0.3pt}
\title{\vspace{-1.3em}\LARGE\bfseries Smooth isometric embeddings of closed\newline orientable surfaces in four-space}
\author{}
\date{9 October 2026\quad\textbullet\quad Public draft 1}

\begin{document}
\maketitle
\thispagestyle{plain}

\section{Problem and notation}
\label{sec:statement}

All manifolds, maps, and metrics are smooth, meaning $C^\infty$.
A closed manifold is compact and has no boundary. Write $\delta_N$ for the
Euclidean metric on $\R^N$ and $\gamma(F)=F^*\delta_N$ for the metric induced
by an immersion. Fix a smooth background metric $g_{\mathrm{bg}}$ and a finite
atlas when taking norms. All $o(1)$ estimates below are uniform on the
specified closed sets.

\begin{theorem}[Smooth isometric embedding]
\label{thm:main}
Every closed orientable
smooth surface $M$ with any smooth Riemannian metric $g$ admits a smooth
embedding $f:M\hookrightarrow\R^4$ satisfying
\begin{equation}
 \langle \dd f_p(v),\dd f_p(w)\rangle=g_p(v,w)
 \qquad(p\in M,\ v,w\in T_pM).
 \label{eq:main}
\end{equation}
No curvature condition is imposed.
\end{theorem}

Equation~\eqref{eq:main} expresses preservation of the Riemannian metric
on tangent spaces. The proof uses exact primitive realization,
boundary propagation, and uniform spherical preparation in the forms
stated below. We show that these operations preserve global injectivity.

\section{Primitive metric corrections}
\label{sec:inputs}

We use exact primitive addition with the following local geometric data
and boundary conditions. Our convention for the vector-valued second
fundamental form is
\begin{equation}
 F_{ij}=\Gamma^k_{ij}F_k+B_F(\partial_i,\partial_j).
 \label{eq:gauss-decomposition}
\end{equation}

\begin{definition}[Admissible primitive data]
\label{def:admissible}
Let $F:M\to\R^4$ be an immersion of a closed surface, with $B_F$ nowhere
the zero tensor and with a smooth global unit normal $n$.
Let $D$ be an open coordinate disk with smooth boundary $C_0$, and let
$(x,y)$ be coordinates on a neighborhood of $\overline D$.
Take a smooth function $a\ge0$ on $M$, positive exactly on $D$, and put
$Y_0=\partial_y$. Require
\begin{equation}
 \Hess_{\gamma(F)}x(\partial_y,\partial_y)>0
 \quad\text{on }\overline D.
 \label{eq:hessian}
\end{equation}
Let $C_1,\ldots,C_m$ be later smooth embedded closed curves, with smooth
nonzero vector fields $Y_i$ on their neighborhoods. The fields need not
be tangent to the curves. Require:
\begin{enumerate}[label=(\roman*)]
\item The curves $C_0,\ldots,C_m$ are pairwise transverse and have no
triple intersection.
\item On $C_i\cap\overline D$, for $i>0$, the field $\partial_y$ is
tangent to $C_i$ at only finitely many points.
\item At $C_i\cap C_l\cap D$, for distinct $i,l>0$, both $\dd x(Y_i)$
and $\dd x(Y_l)$ are nonzero.
\item With $A_i=B_F(Y_i,Y_i)$ and
\begin{equation}
 Q_{il}=\left\langle B_F(Y_i,Y_l),\frac{A_i}{|A_i|}\right\rangle^2
       +K_{\gamma(F)}|Y_i\wedge Y_l|_{\gamma(F)}^2,
 \label{eq:Q}
\end{equation}
the following conditions hold for all indices $0,\ldots,m$:
\begin{equation}
 \begin{aligned}
 A_i&\ne0,\qquad A_i/|A_i|\ne-n &&\text{on }C_i,\\
 \langle A_i,n\rangle&>0,\qquad Q_{il}>0
             &&\text{on }C_i\cap C_l\quad(i\ne l).
 \end{aligned}
 \label{eq:boundary}
\end{equation}
The crossing conditions are imposed in both orders.
\end{enumerate}
\end{definition}

\subsection{Exact realization and uniform asymptotics}
For these data, the velocity-loop construction uses the plane, projected
velocity, and radius
\begin{equation}
 P_p=\Span(F_y(p),F_{yy}(p))^\perp,\qquad
 v=\pi_PF_x,\qquad R_0=(|v|^2+a^2)^{1/2}.
 \label{eq:plane}
\end{equation}
There is a smooth $2\pi$-periodic velocity $V(p,\theta)$, also smooth in
the two-jet of a nearby map, such that
\begin{equation}
 V\in P,\qquad |V|=R_0>0,\qquad \avg_\theta V=v.
 \label{eq:velocity}
\end{equation}
The difference $V-v$ extends smoothly by zero across $\partial D$.
All loop choices, including the finite turn threshold used to preserve
later boundary data, are fixed before choosing the small scale $z$.

The exact primitive-realization theorem supplies, for every sufficiently
small $z>0$, a smooth immersion $f_z$ such that
\begin{equation}
 \gamma(f_z)=h_*:=\gamma(F)+a^2\dd x^2.
 \label{eq:exact}
\end{equation}
The primitive tensor is extended by zero. With $X=F_x-v+V$, the same
theorem gives the uniform estimates on $\overline D$
\begin{equation}
 \begin{aligned}
 (f_z)_x&=F_x-v+V(p,x/z)+o(1),& (f_z)_y&=F_y+o(1),\\
 (f_z)_{yy}&=F_{yy}+o(1),& (f_z)_{xy}&=X_y+o(1),\\
 z(f_z)_{xx}&=V_\theta+o(1).
 \end{aligned}
 \label{eq:asymptotics}
\end{equation}
Periodic arguments are evaluated at $\theta=x/z$; $X_y$ holds $\theta$
fixed. On $M\setminus D$, one has $f_z\to F$ in $C^2$.

These estimates hold for the \emph{final exact maps}: the final correction
is $o(1)$ in $C^2$ and retains the displayed asymptotics. The realization
theorem also supplies a uniform positive lower bound for the size of the
second fundamental form.

\subsection{The same maps preserve the later conditions}
The boundary-propagation theorem uses the same exact maps with the loop
threshold fixed first. For all sufficiently small $z$, it supplies a global
preferred unit normal $n_z$, preserves the nowhere-zero second-form
tensor, and establishes \eqref{eq:boundary} for all later indices.
The adjustments at crossings and at the collar are adjustments of
normal fields; they do not replace $f_z$. Consequently an additional
condition that holds for all sufficiently small $z$ may be imposed on
these same exact maps.

\section{A geometric stability lemma}
\label{sec:stability}

The results in this section do not depend on the analytic construction.
Let $F:M\hookrightarrow\R^N$ be an embedding of a closed manifold, and
write $T_p=\dd F_p(T_pM)$.

\begin{lemma}[Tubular coordinates for a complement]
\label{lem:tubular}
Let $E$ be a smooth subbundle of $M\times\R^N$ such that
$\R^N=T_p\oplus E_p$ for every $p$. For some $r>0$, the map
\[
 \Psi:E_{<r}\longrightarrow\R^N,\qquad
 \Psi(p,e)=F(p)+e,\qquad E_{<r}=\{(p,e)\in E:|e|<r\},
\]
is a diffeomorphism onto an open neighborhood $U$ of $F(M)$.
The retraction $H=F\circ\pi\circ\Psi^{-1}:U\to F(M)$ satisfies
\begin{equation}
 DH_{F(p)}|_{E_p}=0,\qquad DH_{F(p)}\dd F_p=\dd F_p.
 \label{eq:retraction}
\end{equation}
\end{lemma}
\begin{proof}
At the zero section, the natural tangent splitting gives
$D\Psi_{(p,0)}(v,e)=\dd F_p(v)+e$, an isomorphism. The inverse
function theorem and compactness imply that $\Psi$ is a local
diffeomorphism on $E_{<r}$ for a sufficiently small uniform $r$.

If no smaller radius gave injectivity, there would be unequal pairs
$(p_j,e_j),(q_j,e'_j)$ with $|e_j|,|e'_j|\to0$ and
$F(p_j)+e_j=F(q_j)+e'_j$. After taking subsequences, $p_j\to p$ and
$q_j\to q$. Then $F(p)=F(q)$, so $p=q$. Both pairs converge in the
total space to $(p,0)$, contradicting local injectivity there.
Thus $\Psi$ is an injective local diffeomorphism after shrinking $r$.
Its image is open and its inverse smooth. The identities follow by
differentiating $H(F(p)+e)=F(p)$.
\end{proof}

\begin{lemma}[$C^1$ openness of compact embeddings]
\label{lem:open}
If smooth maps $G_j:M\to\R^N$ converge to $F$ in $C^1$, then $G_j$
is an embedding for all sufficiently large $j$.
\end{lemma}
\begin{proof}
Near any point choose coordinates $u$ and a linear map
$L:\R^N\to\R^{\dim M}$ so that $D(LF)$ at the center is the identity.
On a sufficiently small convex coordinate ball, its derivative differs
from the identity by less than $1/4$. For all sufficiently large $j$,
$D(LG_j)$ differs from the identity by less than $1/2$ there. Integrating
the derivative error along a segment yields
\[
 |LG_j(u)-LG_j(v)|\ge\tfrac12|u-v|.
\]
This proves injectivity and immersion on that ball. Take finitely many
such balls whose smaller concentric balls cover $M$. The union of the
smaller-ball products is a neighborhood of the diagonal in $M\times M$.
Its complement is compact and disjoint from the diagonal. If nonempty,
$|F(p)-F(q)|$ has a positive minimum there; uniform convergence preserves
this separation. Thus $G_j$ is an injective immersion on all of $M$.
A continuous injection from compact $M$ into a Hausdorff space is a
homeomorphism onto its image.
\end{proof}

\begin{proposition}[Embedding preservation under transverse oscillations]
\label{prop:stability}
Let $E$ be as in Lemma~\ref{lem:tubular}. Suppose smooth maps
$f_j:M\to\R^N$ satisfy
\begin{equation}
 f_j\to F\text{ in }C^0,\qquad
 \sup_j\|\dd f_j\|_{C^0}<\infty,\qquad
 \dd f_j-\dd F=W_j+R_j,
 \label{eq:abstract-assumptions}
\end{equation}
where $W_{j,p}(T_pM)\subset E_p$ and $\|R_j\|_{C^0}\to0$.
Then $f_j$ is an embedding for all sufficiently large $j$.
\end{proposition}
\begin{proof}
For large $j$, uniform convergence puts $f_j(M)$ in the fixed tubular
neighborhood $U$. At each $p$, the chain rule gives
\begin{align}
 \dd(H\circ f_j)_p-\dd F_p
 &=\bigl(DH_{f_j(p)}-DH_{F(p)}\bigr)\dd f_{j,p}
   +DH_{F(p)}(\dd f_{j,p}-\dd F_p)\notag\\
 &=\bigl(DH_{f_j(p)}-DH_{F(p)}\bigr)\dd f_{j,p}
   +DH_{F(p)}R_{j,p}.
 \label{eq:cancellation}
\end{align}
The term $W_j$ vanishes exactly under $DH_{F(p)}$.
On a compact subneighborhood of $F(M)$, $DH$ is uniformly continuous.
The first term therefore tends uniformly to zero by the boundedness
of $\dd f_j$, and the second does so by the assumption on $R_j$.
Also $H\circ f_j\to F$ uniformly, so the convergence is in $C^1$.
Lemma~\ref{lem:open} makes $H\circ f_j$ an embedding for large $j$.

If $f_j(p)=f_j(q)$, injectivity of the composition gives $p=q$.
If $\dd f_{j,p}(v)=0$, injectivity of $\dd(H\circ f_j)_p$ gives $v=0$.
Thus $f_j$ is an injective immersion, and compactness finishes the proof.
\end{proof}

\begin{remark}
The same argument applies to a family indexed by $z\downarrow0$, with
no continuity in $z$ required. The composition $H\circ f_z$ is only a
test for injectivity; it does not replace the exact isometric map $f_z$.
No metric identity is asserted for the composition.
\end{remark}

\section{Embedding preservation in one primitive addition}
\label{sec:primitive}

\begin{lemma}[Transversality of the oscillation plane]
\label{lem:plane}
Under \eqref{eq:hessian}, the planes $P$ in \eqref{eq:plane} form a
smooth two-plane bundle near $\overline D$, and
\begin{equation}
 \R^4=\dd F_p(T_pM)\oplus P_p.
 \label{eq:complement}
\end{equation}
\end{lemma}
\begin{proof}
Since $x$ is a coordinate function,
$\Hess_{\gamma(F)}x(\partial_y,\partial_y)=-\Gamma^x_{yy}$.
Thus $\Gamma^x_{yy}\ne0$. By \eqref{eq:gauss-decomposition},
\[
 F_{yy}=\Gamma^x_{yy}F_x+\Gamma^y_{yy}F_y+B_{yy}.
\]
The tangential projections of $F_y$ and $F_{yy}$ are independent,
so the ambient vectors are independent and their perpendicular
complement is a smooth two-plane bundle. If $w\in P_p\cap\dd F_p(T_pM)$,
then $\langle w,F_y\rangle=\langle w,F_{yy}\rangle=0$ and
$\langle w,B_{yy}\rangle=0$. Hence
$\Gamma^x_{yy}\langle w,F_x\rangle=0$, which implies $w=0$.
Strict positivity on the compact set $\overline D$ persists on a
neighborhood, so the assertions hold there as well.
\end{proof}

The associated determinant is
\begin{equation}
 \det\begin{pmatrix}
 \langle F_y,F_x\rangle&\langle F_y,F_y\rangle\\
 \langle F_{yy},F_x\rangle&\langle F_{yy},F_y\rangle
 \end{pmatrix}
 =-\Gamma^x_{yy}\det\gamma(F)>0.
 \label{eq:determinant}
\end{equation}
This transversality will let us remove the leading oscillation by
projection along a complement to the old tangent bundle.

\begin{remark}[Nonvanishing of the velocity radius]
The Hessian condition alone does not imply $v\ne0$ everywhere.
In $D$, $a>0$ gives $|v|^2+a^2>0$. On $\partial D$, the boundary
condition gives $B_{yy}\ne0$ and projection yields
$\pi_PB_{yy}=-\Gamma^x_{yy}v$. Projection from the normal plane to $P$
is an isomorphism: its kernel would lie in
$\Span(F_y,F_{yy})$ with zero tangential projection, and that projection
is injective. Hence $v\ne0$ at the boundary and on a collar.
Compactness supplies a uniform positive radius margin.
\end{remark}

\begin{lemma}[Extension of a local complement]
\label{lem:extend}
If $F$ is an embedding, there is a smooth global complement $E$ to
$\dd F(TM)$ agreeing with $P$ near $\overline D$.
\end{lemma}
\begin{proof}
Write $T=\dd F(TM)$ and $N=T^\perp$. On the neighborhood where $P$
is defined, $\pi_N|_P:P\to N$ is a bundle isomorphism. Thus
$P_p=\{A_pu+u:u\in N_p\}$ for a smooth bundle map $A:N\to T$.
Multiply $A$ by a cutoff equal to one near $\overline D$ and supported
in that neighborhood, and extend the product by zero to a global
smooth bundle map $\overline A$. Its graph is the required $E$.
Every such graph is transverse to $T$, regardless of its norm.
\end{proof}

\begin{lemma}[Uniform displacement]
\label{lem:displacement}
The exact maps in \eqref{eq:asymptotics} converge to $F$ in $C^0(M)$.
\end{lemma}
\begin{proof}
Set $e_z=f_z-F$. Exterior convergence controls $e_z$ uniformly on
$\partial D$, while \eqref{eq:asymptotics} gives
$\|\partial_y e_z\|_{C^0(\overline D)}\to0$.
For a point $(x_0,y_0)\in D$, take the connected interval containing
$y_0$ in the vertical slice $\{y:(x_0,y)\in D\}$.
It has an endpoint on $\partial D$ and length at most the finite
$y$-range $L_D$ of the coordinate image of $\overline D$.
Integrating to that endpoint gives
\begin{equation}
 \|e_z\|_{C^0(\overline D)}\le
 \|e_z\|_{C^0(\partial D)}+
 L_D\|\partial_y e_z\|_{C^0(\overline D)}\longrightarrow0.
 \label{eq:displacement}
\end{equation}
Combine this with exterior convergence. Vertical convexity of the
whole disk and continuous choices of endpoints are unnecessary.
\end{proof}

\begin{proposition}[Embedding-preserving primitive addition]
\label{prop:primitive}
For admissible primitive data, suppose that $F$ is an embedding.
The exact primitive output $F_+$ can be chosen to be an embedding,
retaining its target metric,
nowhere-zero second-form tensor, preferred normal, and later boundary
conditions. Its exterior $C^2$ error and global $C^0$ displacement
can be made arbitrarily small.
\end{proposition}
\begin{proof}
Fix the velocity loop, and let $f_z$ be its exact realizations. For the global
complement from Lemma~\ref{lem:extend}, set
\[
 W_z=(V(p,x/z)-v(p))\,\dd x\quad\text{on }D,
 \qquad W_z=0\quad\text{outside }D.
\]
For each fixed $z$, this is a smooth global one-form with values in $E$,
by the smooth zero extension of $V-v$. The first-derivative asymptotics
and exterior convergence give
\[
 \dd f_z-\dd F=W_z+R_z,\qquad \|R_z\|_{C^0(M)}\to0.
\]
For background-unit $u$, the exact metric identity gives
$|\dd f_z(u)|^2=h_*(u,u)$. Since $h_*$ is fixed, the differentials
are uniformly bounded. Lemma~\ref{lem:displacement} gives the remaining
hypothesis of Proposition~\ref{prop:stability}, so $f_z$ is an embedding
for all sufficiently small $z$.

By the boundary-propagation theorem in Section~\ref{sec:inputs}, the other
geometric conclusions hold for these same maps for every sufficiently small
$z$. Choose $z$ below all the thresholds, also meeting the requested
exterior error and displacement, and set $F_+=f_z$.
\end{proof}

\section{Embedded preparation and finite assembly}
\label{sec:global}

\begin{lemma}[Embedded spherical preparation]
\label{lem:preparation}
For a closed connected orientable surface, uniform spherical preparation
and global metric preparation can be performed with the initial map
$F^0$ an embedding, retaining the required geometric conditions and
prescribed intermediate metrics.
\end{lemma}
\begin{proof}
A closed connected orientable smooth surface is diffeomorphic to the
boundary of a handlebody in $\R^3$. Compose such an embedding with
inverse stereographic projection to obtain $I:M\hookrightarrow\Sph^3$.
No metric is prescribed for this starting embedding.

\paragraph{A uniform family inside the embedding neighborhood.}
The uniform spherical-preparation theorem applies to this fixed $I$ and every finite set
$P\subset M$. It supplies constants $C,c,\varepsilon_0>0$ such that,
for any $0<\varepsilon<\varepsilon_0$, one can choose $I_P:M\to\Sph^3$ with
\begin{equation}
 \begin{gathered}
 B^{\Sph^3}_{I_P}|_P=0,\qquad \|I_P-I\|_{C^1}<\varepsilon,\\
 \|I_P\|_{C^2}\le C,\qquad I_P^*\delta_4\ge c\,g_{\mathrm{bg}}.
 \end{gathered}
 \label{eq:spherical-preparation}
\end{equation}
The constants are independent of the number and separation of points
of $P$. Fix $\varepsilon$ small enough for Lemma~\ref{lem:open}.
Then every permitted $I_P$ is an embedding, including maps whose finite
set is chosen later. Restrict the preparation family to this neighborhood.

\paragraph{The order of choices in the global construction.}
Global metric preparation begins with a fixed spherical immersion and
its uniform preparation family. Use the embedding $I$ just chosen.
First choose $\lambda>0$ uniformly for the
restricted family so that
\[
 \gamma^0_{\overline I}=\lambda^2\overline I^*\delta_4<g,
 \qquad h_{\overline I}=g-\gamma^0_{\overline I}>0.
\]
The bounds in \eqref{eq:spherical-preparation} give uniform positive
lower bounds and $C^1$ bounds for all path metrics
$\gamma^0_{\overline I}+s h_{\overline I}$, $0\le s\le1$.
In particular their Christoffel symbols are uniformly bounded.

Uniform metric preparation gives convex phases, disks, and allowed
perturbation sizes for this family, with strict cone and phase-Hessian
margins. Choose a number $N_c$ of cycles sufficiently large to retain
these margins at every partial metric. Each perturbed cycle has an
invertible bundle endomorphism $\mathcal A_c$ of symmetric tensors and
amplitudes
\begin{equation}
 a_j=N_c^{-1/2}\psi_j
       \sqrt{q^0_{\nu r}(\mathcal A_c^{-1}h_{\overline I})},
 \qquad
 \sum_{j\text{ in cycle }c}a_j^2\dd x_j^2=N_c^{-1}h_{\overline I}.
 \label{eq:cycles}
\end{equation}
Here $q^0_{\nu r}$ are the coefficient functionals of its positive
rank-one decomposition, and $\psi_j$ are its support cutoffs.
These formulas and their uniform bounds are part of the metric-preparation
theorem. Only the already fixed bounds on the family enter these choices.

Now choose the allowed generic boundary and phase perturbations.
Their boundary-crossing set $P$ is finite. Choose $I_P$ only at this
stage and set $F^0=\lambda I_P$, $n^0=-I_P$. Insert the actual deficit
$h=g-\gamma(F^0)$ in \eqref{eq:cycles}. The uniform preparation estimates
retain the cone, Hessian, and exact-cycle identities for this choice.
Since $I_P$ is an embedding, so is $F^0$. No earlier constant needs to
be chosen again after selecting $P$.

\paragraph{Initial boundary geometry.}
For $\gamma^0=\gamma(F^0)$, the radial decomposition is
\begin{equation}
 B_{F^0}(v,w)=\lambda B^{\Sph^3}_{I_P}(v,w)
                +\lambda^{-1}\gamma^0(v,w)n^0.
 \label{eq:radial}
\end{equation}
Its dot product with $n^0$ is positive for $v=w\ne0$, so the second-form
tensor is nowhere zero and every pure direction avoids $-n^0$.
At a point of $P$, the spherical term vanishes. The Gauss equation gives
$K_{\gamma^0}=\lambda^{-2}$ and
\begin{equation}
 \begin{aligned}
 Q_{il}(F^0)
 &=\lambda^{-2}\bigl(\gamma^0(Y_i,Y_l)^2+
                         |Y_i\wedge Y_l|_{\gamma^0}^2\bigr)\\
 &=\lambda^{-2}\gamma^0(Y_i,Y_i)\gamma^0(Y_l,Y_l)>0.
 \end{aligned}
 \label{eq:initial-crossing}
\end{equation}
This verifies both ordered crossing conditions. All initial hypotheses
of the finite construction are therefore retained.
\end{proof}

\begin{proof}[Proof of Theorem~\ref{thm:main}]
Assume first that $M$ is connected. Lemma~\ref{lem:preparation} gives
an embedded $F^0$ and a finite list $(D_j,x_j,y_j,a_j)$, $1\le j\le K$,
with exact prescribed metrics
\begin{equation}
 \begin{gathered}
 \gamma^j=\gamma(F^0)+\sum_{i=1}^j a_i^2\dd x_i^2,
 \qquad \gamma^K=g,\\
 \Hess_{\gamma^{j-1}}x_j(\partial_{y_j},\partial_{y_j})>0
 \quad\text{on }\overline D_j.
 \end{gathered}
 \label{eq:finite-metrics}
\end{equation}
Suppose $F^{j-1}$ is an embedding with induced metric $\gamma^{j-1}$,
nowhere-zero second-form tensor, and a preferred normal satisfying all
remaining boundary conditions. The geometric curve conditions are fixed
by preparation; the Hessian condition is intrinsic and follows from
\eqref{eq:finite-metrics}. Proposition~\ref{prop:primitive} gives an
embedding $F^j$ inducing $\gamma^j$ and retaining the conditions needed
for the next step. This closes the finite induction. Its final map
$f=F^K$ is an embedding with induced metric $g$.

Each scale is chosen after fixing the current smooth map. Constants
and tubular radii need not be uniform between stages. At each stage,
Proposition~\ref{prop:primitive} applies to the exact output of the
primitive-realization theorem. No
injectivity assertion is inferred from a mere $C^0$ limit of embeddings.

A compact manifold has finitely many connected components because
its components are open. Apply the connected construction to each
component and translate the finitely many compact images to make them
disjoint. Translations preserve the induced metric. The empty case is
vacuous.
\end{proof}

\section{Conclusion}
\label{sec:conclusion}

The construction starts from an embedded surface in a small round
three-sphere and decomposes the prescribed metric change into finitely many
primitive additions. At each step, the positive transverse phase Hessian
makes the oscillation plane complementary to the old tangent plane.
An oblique tubular retraction cancels the leading derivative oscillation,
so its composition with the exact corrected map is $C^1$ close to the
previous embedding. This proves injectivity of the corrected map itself.
The corrected map retains the exact new metric and the geometric
conditions needed for the next addition.

Uniform spherical preparation provides the initial embedding and the
finite metric sequence. Applying the embedding-preserving primitive step
along this sequence gives a final smooth embedding with the prescribed
metric. Thus, using the exact primitive-realization, boundary-propagation,
and preparation theorems in the forms stated above, we conclude that every
closed orientable smooth Riemannian surface $(M,g)$ admits a smooth embedding
\[
 \boxed{f:M\hookrightarrow\R^4,\qquad f^*\delta_4=g.}
\]
The metric $g$ is arbitrary; no curvature restriction is required. For disconnected surfaces, translations separate the finitely
many component images while preserving their induced metrics.

\end{document}
